% Compile with XeLaTeX twice. Official result fields must be filled from actual runs.
\documentclass[UTF8,a4paper,zihao=-4,fontset=windows]{ctexart}
\usepackage[margin=2.45cm]{geometry}
\usepackage{amsmath,amssymb,amsthm,booktabs,longtable,array,graphicx}
\usepackage{tikz}
\usepackage[hidelinks]{hyperref}
\usetikzlibrary{arrows.meta}
\linespread{1.17}
\setlength{\parindent}{2em}
\setlength{\emergencystretch}{2em}
\newtheorem{theorem}{定理}[section]
\newtheorem{proposition}[theorem]{命题}
\newcommand{\norm}[1]{\left\lVert#1\right\rVert}
\newcommand{\R}{\mathbb R}
\newcommand{\pending}{\textbf{待实测}}
\newcommand{\uvec}{u}
\title{有界观测下无线电干扰源的\\覆盖证书与搜索清除联合调度}
\author{}
\date{}
\begin{document}
\maketitle
\vspace{-2.8em}
\begin{center}\small 第七版正文稿：纳入第三、六版真实演练；第七版实测及正式结果待补\end{center}

\begin{abstract}
针对源总数未知、接收半径未知、测角误差有界且辐射方向不完全可见的干扰源清除问题，本文联合维护全局搜索覆盖、历史相容区域和可执行清除策略，不依赖误差独立性或源位置概率分布。

问题一给出前向楔形交直径算法及同直径圆不能覆盖定位区域的可实现反例。问题二消去固定未知接收半径，以解析直径界不超过163米为约束，同时优化移动距离与侧向角度，得到886.3434米可行移动距离，距该解析判据最优值的计算间隙不超过0.0820米；该结论不等于真实物理直径的精确极小极大最优性。

问题三采用原点与六点环的七点覆盖，问题四采用原点、8个内环点和16个外环点的25点三角网，以连续区域证书保证发现。在执行层，将全部待处理源与剩余扫描统一为有向服务图：以实际第一测点作为入口、整个可能清除区域作为出口，结合指派松弛、环连接和固定顺序下的模式动态规划重排全部任务。每轮执行一个任务后更新，允许连续清除多个源。

对整个相容区域求最近的保证收缩测点，代替只按旧锚点距离折半；定向源仅在所有相容发射方向均可见时采用单点替代。额外观测用完整示向区间评价；光学终局按此前失败条件剪除不可能成功的尝试，仍以光学上界严格低于测向策略下界作为切换门槛。覆盖改写分阶段保留证明预算。上述替代保留五次全向或十二次混合核心测量保证，以及300、772次动作上界。

第七版12项必要固定输入检查通过。第六版真实演练分别清除13、14源，虚拟时间4889.66、8753.57秒，移动占85.07\%、83.78\%；两例官方总源数未提供。真实首批观测上的计划界分析支持新的调度方向，但不构成第七版实测提速。第七版实际性能与六次正式结果仍待操作者运行取得。
\end{abstract}

\noindent\textbf{关键词：}覆盖证书；环形扇区；未知定向；服务图；区域探测
\newpage

\section{问题重述与分析}
目标区域为以原点为中心、半径1800米的圆域。10至16个静止干扰源分别占用20个频道中的不同频道，单源有效接收半径在1000至1500米之间。全向源向全部方向辐射，定向源仅在某未知闭半平面内辐射。机器狗以5米每秒移动，测向需5秒，切换测向频道需1秒，20米以内可通过光学精确定位并清除，总耗时5秒\cite{problem,manual,protocol}。

问题一要求利用若干示向度计算交会定位区域直径，并判断相同直径的圆是否足以覆盖该区域。问题二要求仅依据一次有效示向度选择第二检测点，兼顾可接收性和定位几何条件。问题三、四要求在源总数未知时清除全部源，并分别处理全向与混合辐射情形。

本题存在三项核心困难。首先，同一地点误差固定，重复检测不能通过平均必然提高精度；其次，无信号同时可能表示无源、超距或方向不可见；最后，定位区域很小与机器狗已经到达清除范围并不等价，单纯优化位置估计精度可能增加移动成本。

据此，本文采用分层决策：先构造覆盖证明确保不漏源；再用已知可见点建立局部不变量，主动排除无信号的距离歧义；最后直接寻找可清除位置。问题一的纯交会区域、问题二含距离约束的相容集合以及后两问用于行动的外包区域分别定义，避免混用其几何性质。

\section{题设条件、建模假设与符号}
\subsection{采用的条件与补充约定}
源位置、频道、有效接收半径与发射方向在一次任务中固定，不同源互不影响；移动空间无遮挡且允许离开源分布圆域。这些条件来自题面和附件。算法仅使用检测与清除响应，不读取隐藏源数据。

将每次方向误差视为任意有界扰动，不额外假设独立、正态或均匀分布。纯几何分析采用$\varepsilon=1^\circ$；对应接口两位小数舍入，程序采用$\varepsilon=1.005^\circ$。本文除特别说明外均使用后一个保守值，内部三角函数按弧度计算。

设置安全清除半径$r_s=19.8$米，低于题设20米。证明首先在精确算术下给出；浮点实现采用向外放宽半平面、向内缩小行动半径的余量，并对不相容结果停止处理。这不等价于严格区间算术验证。

\begin{longtable}{p{.22\textwidth}p{.69\textwidth}}
\caption{主要符号}\label{tab:symbols}\\
\toprule 符号 & 含义\\\midrule\endfirsthead
\toprule 符号 & 含义\\\midrule\endhead
$x_c,R_c,n_c$ & 频道$c$源的位置、接收半径、定向半平面的单位法向量\\
$p,h$ & 机器狗当前位置、当前测向频道\\
$u(\theta),v(\theta)$ & $(\cos\theta,\sin\theta)$及其逆时针垂直单位向量\\
$W_i,P$ & 单次前向测角楔形、纯交会定位区域\\
$(a_c,\theta_c,L_c)$ & 可见锚点、锚点示向度、锚点到源的距离上界\\
$\widehat P_c$ & 包含真实源的历史约束外包多边形\\
$Q_3,Q_4$ & 全向、混合情形的固定搜索点集\\
$U,D,C$ & 未发现频道集、全部已发现频道集、成功清除频道集\\
$T_v,T_r$ & 虚拟任务时间、官方程序运行时间\\
\bottomrule
\end{longtable}

\subsection{统一观测与目标函数}
对定向源，检测位置$q$能够收到信号的条件为
\begin{equation}
\norm{q-x_c}\le R_c,\qquad n_c^\mathsf T(q-x_c)\ge0.
\end{equation}
全向源只要求第一个条件。可接收且距离不超过5米时返回\texttt{near}，否则返回带误差的方向；不可接收时返回\texttt{no\_signal}。

记相邻动作的位置为$p_{j-1},p_j$，动作总数为$A$，测量次数为$M$，成功和失败清除次数为$C_s,C_f$，实际测向频道切换次数为$S$。虚拟总时间为
\begin{equation}\label{eq:cost}
T_v=\frac15\sum_{j=1}^{A}\norm{p_j-p_{j-1}}+5M+S+5C_s+3C_f.
\end{equation}
\texttt{/clear}中的频道仅标识清除目标，不改变测向频道。直接清除须有距离保证；光学覆盖终局允许前几次清除失败，失败响应消除对应20米圆内的可能位置，因此$C_f$不再恒为零。目标是在全部清除的约束下减小$T_v$，而非单独最小化测量次数；所给在线调度是确定启发式，不声称求得整局全局最优。

\section{问题一：交会区域与直径}
\subsection{前向楔形的半平面表示}
检测位置$p_i$的示向度为$\theta_i$。用二维叉积$[z,w]=z_xw_y-z_yw_x$定义
\begin{equation}
W_i=\{x:[u(\theta_i-\varepsilon),x-p_i]\ge0,
\ [u(\theta_i+\varepsilon),x-p_i]\le0\},\quad P=\bigcap_iW_i.
\end{equation}
由于$2\varepsilon<180^\circ$，这是前向凸楔形，不能用无限直线两侧的带状区域替代。每个楔形提供两个线性不等式$g_j^\mathsf Tx\le b_j$，因此$P$是凸多面集。

\subsection{分类与可实现算法}
若没有观测，则$P=\R^2$。至少一个尖楔形使交集不含整条直线，若交集非空则存在极点。因此可枚举非平行边界线对的交点，保留满足全部半平面约束的点。若没有可行交点，精确算术下为空集；数值计算中还可能是近平行交点无法解析，程序保守返回“空或数值未决”。

在存在可行点的前提下，考察衰退锥
\begin{equation}
\mathcal R=\{d:g_j^\mathsf Td\le0,\ \forall j\}.
\end{equation}
若存在非零$d$，则$P$无界、直径无穷；否则有界。在二维尖锥条件下，非零衰退锥必有边界射线，可枚举各约束法向量的正负垂直方向，检验全部齐次不等式。

设观测数为$n$。枚举$O(n^2)$个边界交点，每点检查$2n$条约束，总代价$O(n^3)$。对可行顶点计算凸包得到逆时针有序序列$z_1,\ldots,z_m$，随后用旋转卡壳在$O(m)$时间内计算最远点对。该实现适合小规模测向数据，本文不把边界枚举误称为$O(n\log n)$半平面交算法。

\begin{proposition}
若$P$是非空有界凸多边形，则
\begin{equation}\operatorname{diam}(P)=\max_{1\le i,j\le m}\norm{z_i-z_j}.\end{equation}
\end{proposition}
\begin{proof}
固定$y\in P$，若$x=\sum_i\lambda_i z_i$，则$\norm{x-y}\le\sum_i\lambda_i\norm{z_i-y}\le\max_i\norm{z_i-y}$。再对$y$作同样展开即可。单点和线段分别返回0和端点距离。
\end{proof}

\subsection{相同直径的圆并不一定覆盖定位区域}
取边长为$d$的正三角形。其区域直径为$d$，但最小包围圆半径为$d/\sqrt3>d/2$，所以不存在直径为$d$的覆盖圆。对任一最远点对$z_i,z_j$，令$c=(z_i+z_j)/2$，则该直径圆覆盖$P$的充要条件为
\begin{equation}\label{eq:diamcover}
\max_k\norm{z_k-c}\le d/2.
\end{equation}
任何半径$d/2$且同时覆盖该最远点对的圆都必须以$c$为圆心，因此式\eqref{eq:diamcover}也回答了是否存在这样的圆。取纯几何误差半宽为1度时，该反例也可由本题楔形实际构成：沿单位正三角形各条逆时针边反向延长100米放置检测点，示向度取对应边方向加$1^\circ$。每个楔形的一条边界恰为三角形边，另一条边界在三角形外构成冗余约束，三个楔形之交正是该三角形。

\begin{figure}[htbp]\centering
\begin{tikzpicture}[scale=2]
\draw[thick] (-.5,0)--(.5,0)--(0,.866)--cycle;
\draw[dashed] (0,0) circle (.5);
\draw[blue] (0,.288675) circle (.57735);
\fill (-.5,0) circle (.02) node[below left] {$A$};
\fill (.5,0) circle (.02) node[below right] {$B$};
\fill (0,.866) circle (.02) node[above] {$C$};
\node at (1.35,.65) {虚线圆：$r=d/2$};
\node[blue] at (1.5,.3) {包围圆：$r=d/\sqrt3$};
\end{tikzpicture}
\caption{相同直径圆不能覆盖区域的解析反例}
\end{figure}

\section{问题二：保证接收的第二测点与交会精度}

\subsection{利用首次接收消去未知半径}
令第一次正常测向位置为$p_1$，示向度为$\theta_1$，取$u=u(\theta_1),v=v(\theta_1)$，记$L=1500$米。真实源满足
\begin{equation}
x=p_1+r(\cos\delta\,u+\sin\delta\,v),\quad 5<r\le1500,\quad |\delta|\le\varepsilon.
\end{equation}
第一次已经收到信号，故$R\ge\max\{1000,r\}$。用$r\in[0,1500]$放宽先验，设第二点$q=p_1+au+bv$、$a\ge0$，要求
\begin{equation}\label{eq:coupling}
\norm{q-p_1-r u(\theta_1+\delta)}\le\max\{1000,r\},\quad\forall r,\delta.
\end{equation}
当$0\le r\le1000$时，距离平方关于$r$的凸性表明只需检查$r=0,1000$。当$r\ge1000$时，平方消去$r^2$得到
\begin{equation}
a^2+b^2\le2r(a\cos\delta+b\sin\delta).
\end{equation}
只要$r=1000$的条件满足，右侧系数非负，随$r$增大条件不会变差。因此在前向半平面内，放宽先验对应的完整保证接收集合为
\begin{equation}\label{eq:candidate}
\begin{aligned}
\mathcal C_{\rm new}={}&B((0,0),1000)\cap\{a\ge0\}\\
&\cap B((1000\cos\varepsilon,1000\sin\varepsilon),1000)\\
&\cap B((1000\cos\varepsilon,-1000\sin\varepsilon),1000).
\end{aligned}
\end{equation}
对真实的$r>5$及源圆域先验，它仍是保守的充分集合。原先若要求全部1500米方向端点都在1000米内，得到$\mathcal C_{\rm old}$；由径向凸性有$\mathcal C_{\rm old}\subset\mathcal C_{\rm new}$。

例如$(a,b)=(500,840)$距原点约977.55米，距新区域最坏方向端点约992.58米，属于新区域；距旧1500米方向端点约1322.89米，不属于旧区域。这证明扩大严格成立，但不说明该点的定位精度优于原推荐点。
\begin{figure}[htbp]\centering
\begin{tikzpicture}[scale=3]
\begin{scope}
\clip (0,0) circle (1);
\clip (.999846,.01754) circle (1);
\clip (.999846,-.01754) circle (1);
\fill[blue!15] (0,-1.1) rectangle (1.1,1.1);
\end{scope}
\begin{scope}
\clip (0,0) circle (1);
\clip (1.499769,.02631) circle (1);
\clip (1.499769,-.02631) circle (1);
\fill[green!40] (0,-1.1) rectangle (1.1,1.1);
\end{scope}
\draw[->] (-.1,0)--(1.25,0) node[right] {$a/1000$};
\draw[->] (0,-.95)--(0,1.03) node[above] {$b/1000$};
\fill[red] (.5,.84) circle (.014) node[above right] {$(500,840)$};
\fill (.75,.6) circle (.014) node[right] {旧点};
\fill[orange] (.664633,.606847) circle (.014) node[above left] {新点};
\node[blue!70!black,anchor=west] at (1.1,-.3) {蓝色：扩大候选域};
\node[green!40!black,anchor=west] at (1.1,-.5) {绿色：原候选域};
\end{tikzpicture}
\caption{半径消元严格扩大保证接收区域，图中纵横坐标均已归一化}
\end{figure}


\subsection{移动距离约束下的测点优化}
令$q-p_1=\rho(\cos\psi\,u+\sin\psi\,v)$，只优化正侧，负侧由对称性取得。固定从首次测点移动的距离$5<\rho<1000$，式\eqref{eq:candidate}等价于
\begin{equation}
0\le\psi\le\arccos(\rho/2000)-\varepsilon.
\end{equation}
有价值的交会设计还要求$\psi>\varepsilon$。定义
\begin{align}
D(\rho,\psi)&=\max_{r\in\{5,1500\}}
\sqrt{r^2+\rho^2-2r\rho\cos(\psi+\varepsilon)},\\
h(\rho,\psi)&=\rho\sin(\psi-\varepsilon),\\
\gamma_0&=\arcsin(h/D)-2\varepsilon.
\end{align}
径向凸性及最坏角端点给出所有真实源到第二点的距离不超过$D$。真实两条视线的无向锐夹角至少为$\arcsin(h/D)$，故两条示向直线的锐夹角至少为$\gamma_0$。若$\gamma_0\le0$，该设计不提供下面的有限界。

令$w_1=1500\sin\varepsilon,w_2=D\sin\varepsilon$。物理相容集合落在两条宽度分别为$2w_1,2w_2$的带内。计算带交平行四边形的最长对角线，得到比直接三角不等式更紧的界
\begin{equation}\label{eq:163}
\operatorname{diam}(P_{\rm physical})\le
B(\rho,\psi)=
\frac{2\sqrt{w_1^2+w_2^2+2w_1w_2\cos\gamma_0}}{\sin\gamma_0}.
\end{equation}
固定移动距离时，在完整可行角区间内最小化解析上界$B$；第七版进一步以$B\le163$为约束最小化$\rho$。两种问题均以解析带交上界为目标，不等同于真实相容直径的精确极小极大问题。

\paragraph{连续区间分支定界。}
对角区间$[\ell,h]$，$D$及交会高度$\rho\sin(\psi-\varepsilon)$均随$\psi$增加；式\eqref{eq:163}随$D$增加、随交会高度增加而减小。用$D(\rho,\ell)$和$\rho\sin(h-\varepsilon)$代入，得到该区间目标函数的下界；将比值截断至1只会进一步放宽下界。端点与二分中点给出可行设计上界。程序优先细分下界最小的区间，最多1024次二分子区间的中点目标计算，另有3次初始候选计算，保留未处理区间的最小下界，因此预算耗尽时仍报告明确的剩余优化间隙。

\begin{table}[htbp]\centering\small
\caption{固定移动距离下的解析设计，非模拟实验}
\begin{tabular}{rrrrr}\toprule
移动距离/米&前进/米&侧移/米&直径上界/米&优化间隙上限/米\\\midrule
600.00&315.50&510.35&271.673&0.020\\
800.00&539.46&590.74&187.698&0.021\\
900.00&664.63&606.85&159.499&0.023\\
960.47&743.30&608.27&145.165&0.020\\
999.00&763.89&643.80&137.238&0.020\\\bottomrule
\end{tabular}
\end{table}
表中坐标仅供展示，程序使用未舍入坐标。优化间隙指同一固定移动距离下解析目标的上、下界之差，浮点计算并非严格区间算术。

\paragraph{精度约束下的最短移动。}
令目标直径为$b_*=163$米，求解
\[
\min \rho\quad\text{s.t.}\quad B(\rho,\psi)\le b_*,
\quad 5<\rho<1000,\quad
\varepsilon<\psi\le\arccos(\rho/2000)-\varepsilon .
\]
对矩形节点$[\rho_0,\rho_1]\times[\psi_0,\psi_1]$，令
\[
\underline D=\max_{r\in\{5,1500\}}\min_{\rho\in[\rho_0,\rho_1]}
\sqrt{r^2+\rho^2-2r\rho\cos(\psi_0+\varepsilon)},\qquad
\overline h=\rho_1\sin(\psi_1-\varepsilon).
\]
内层一维二次式的最小点是$r\cos(\psi_0+\varepsilon)$向区间的投影。将$\underline D,\overline h$代入带交公式，得到节点内$B$的下界；若其超过目标，整节点不可能提供可行改进。可行候选更新移动距离上界，其余节点以$\rho_0$排序二分，保留全部未处理节点的最小$\rho_0$。节点配额6000，移动间隙停止阈值0.1米。

本次解析计算处理3065节点，得到
\begin{equation}\label{eq:q2}
q=p_1+651.5212558u\pm600.9364774v,
\qquad \rho=886.3433853\ {\rm m},
\end{equation}
其$B=162.9997924$米，三圆保证接收域的最小余量113.6566米。在上述解析优化域内，最优移动距离的计算区间为
\[
886.2614456\le \rho_*\le886.3433854,\qquad
\rho_{\rm feasible}-\rho_*\le0.081940\ {\rm m}.
\]
这是带交解析判据的分支定界结果，浮点实现尚非严格区间算术认证。接收条件外的设计、真实物理直径极小极大以及其他测点政策不在最优性范围内。

较第六版900米默认移动再减少13.6566米；较旧$(750,\pm600)$设计减少74.1253米，折合14.8251秒纯移动，但正常响应的解析精度保证放宽至163米。精度与移动的权衡被明确写入目标，不能把较紧旧界误当作测点本身的精度收益。显示坐标经过舍入，程序使用未舍入结果。有当前位置时从对称两点选更近者；显式给定移动距离仍可调用固定半径角度优化。

\section{共同核心：正常响应下界与仿射收缩}
\subsection{环形扇区不变量}
首次正常方向响应不仅说明可接收，还排除了5米以内位置。因此保存$\mathcal A=(a,\theta,L)$，初始$L=1500$，并保持
\begin{equation}\label{eq:invariant}
x=a+r(\cos\delta\,u+\sin\delta\,v),\quad 5<r\le L,\quad |\delta|\le\varepsilon,
\quad a\text{可见}.
\end{equation}
当$L>44.5$米时令$A=(L+5)/2,B=0.02L$，构造
\begin{equation}\label{eq:pair}
q_\pm=a+Au\pm Bv.
\end{equation}
先测离当前位置较近的一点，等距取正侧。把中心从$[0,L]$中点改为$[5,L]$中点，是利用响应语义增强几何界，而不是引入新的统计假设。

\begin{proposition}[阳性仿射收缩]\label{prop:range}
对式\eqref{eq:invariant}的全部相容源，当$L>44.5$时有
\begin{equation}\norm{q_\pm-x}<0.502L-2.4\le750.6<1000.\end{equation}
\end{proposition}
\begin{proof}
取最坏角误差后，距离平方关于$r$为凸函数，只需检查$r=5,L$；后一个端点的上界更大。因此
\begin{align}
\norm{q_\pm-x}^2
&\le (L-5)^2/4+0.0004L^2+L(L+5)(1-\cos\varepsilon)+0.04L^2\sin\varepsilon\\
&<0.251256L^2-2.49923L+6.25.
\end{align}
其中使用$\cos\varepsilon>0.999846$、$\sin\varepsilon<0.01755$。后一表达式与$(0.502L-2.4)^2$的差的相反数为
\begin{equation}0.000748L^2+0.08963L-0.49>0,\qquad L>44.5.\end{equation}
故结论成立。两个测点均排除了接收距离不足这一无信号原因。
\end{proof}

\begin{proposition}[双阴性上界]\label{prop:negative}
若两个测点均无信号，则保留旧锚点并可取
\begin{equation}L'=0.501L+2.505.\end{equation}
\end{proposition}
\begin{proof}
令$x-a=tu+wv$，则$|w|\le t\tan\varepsilon$。由命题\ref{prop:range}，两点都在接收距离内，双阴性只能源于定向不可见。假设$t>A$，取
\begin{equation}
\lambda_++\lambda_-=\frac{t}{t-A},\qquad
\lambda_+-\lambda_-=\frac{Aw}{B(t-A)}.
\end{equation}
在$L>44.5$时$B>A\tan\varepsilon$，故两个系数非负，且
$a-x=\lambda_+(q_+-x)+\lambda_-(q_--x)$。
两个不可见向量的非负组合不可能成为可见锚点向量，产生矛盾。若$t=A$，源在两测点连线内部，任一闭半平面至少包含一个端点，也矛盾。因此$t<A$，从而
\begin{equation}
r<\frac{L+5}{2\cos\varepsilon}<0.501L+2.505.
\end{equation}
注意此分支不能误用阳性的$0.502L-2.4$。
\end{proof}

\subsection{环形扇区上的光学终止}
令$m=a+(L+5)u/2$，则
\begin{equation}\label{eq:annularradius}
\norm{x-m}\le g(L),\quad
g(L)=\sqrt{(L-5)^2/4+L(L+5)(1-\cos\varepsilon)}.
\end{equation}
$g$在$L\ge5$单调递增，且$g(44.5)\approx19.758577<19.8$。因此$L\le44.5$时无需再次测向。令$s=r_s-g(L)$，清除圆盘$B(m,s)$中所有点均有效，最近候选为
\begin{equation}\label{eq:terminal}
z=\begin{cases}p,&\norm{p-m}\le s,\\m+s(p-m)/\norm{p-m},&\text{否则}.\end{cases}
\end{equation}
程序在$s$上再内缩$10^{-7}$米，减少边界舍入影响。

\begin{theorem}[五轮全向与六轮混合清除]\label{thm:finite}
仅执行本节连续局部收缩时，首次正常发现后，全向源至多新增5次测量，混合情形每源至多新增12次测量，随后一次有距离保证的清除。第七版另外保留每源最多两次额外信息测向配额，这两次不计入本定理。
\end{theorem}
\begin{proof}
正常方向分支换锚点，取$L'=0.502L-2.4$，新正常响应再次保证$r>5$；双阴性保留原锚点及$r>5$，取$L'=0.501L+2.505$。遇到\texttt{near}则立即同点清除。

全向源每轮第一点必可接收，其五次上界依次约为
\begin{equation}
750.6,\quad374.4012,\quad185.5494024,\quad90.7458000,\quad43.1543916.
\end{equation}
第五次后满足终止阈值。混合情形所有分支统一满足$L'\le0.502L+2.505$，从1500出发六轮后不超过28.955275米，也满足阈值；每轮至多两个测点。
\end{proof}

\begin{figure}[htbp]\centering
\begin{tikzpicture}[x=1cm,y=1cm,>=Stealth]
\fill[blue!8] (.5,-.035)--(7,-.5)--(7,.5)--(.5,.035)--cycle;
\draw[->] (0,0)--(7.5,0) node[right] {$u$};
\draw[dashed] (.5,-.35)--(.5,.35) node[above] {$5$};
\draw[dashed] (7,-.7)--(7,.7) node[above] {$L$};
\draw[thick] (3.75,-.85)--(3.75,.85);
\fill (0,0) circle (2pt) node[below] {$a$};
\fill (3.75,.85) circle (2pt) node[above] {$q_+$};
\fill (3.75,-.85) circle (2pt) node[below] {$q_-$};
\node at (3.75,-1.35) {$\text{前进 }(L+5)/2,\quad\text{侧移 }\pm0.02L$};
\end{tikzpicture}
\caption{利用5米内排除信息选择环形扇区中点，图中尺度为示意}
\end{figure}



\section{历史相容区域与清除动作}
\subsection{显式利用接收半径下界}
对已知全向频道，记清除前阳性位置集为$P^+$、阴性位置集为$P^-$，$D=B(0,1800)$。忽略测角前，源位置$x$相容的充要条件是存在$R\in[1000,1500]$使
\begin{equation}
\max\{1000,\max_{p\in P^+}\norm{x-p}\}\le R
<\min_{q\in P^-}\norm{x-q}.
\end{equation}
无阴性时右侧取无穷。因此可以消去$R$，得到
\begin{equation}\label{eq:exactomni}
F_{\rm omni}=D\cap\bigcap_{p\in P^+}B(p,1500)
\cap\bigcap_{q\in P^-}\{x:\norm{x-q}>1000\}
\cap\bigcap_{\substack{p\in P^+\\q\in P^-}}
\{x:\norm{x-p}<\norm{x-q}\}.
\end{equation}
最后一个条件平方相减即为线性不等式
\begin{equation}
2(q-p)^\mathsf T x<\norm q^2-\norm p^2.
\end{equation}

再与历次方向楔形及正常响应下界相交。程序以闭半平面放宽严格不等式，并用32边切线多边形外包阳性圆；源域采用128边切线外包，最大径向过量小于0.543米。新阳性提高同一源的接收半径下界，因此程序会重新处理该频道的旧阴性，而不只使用最新观测。

第七版不再把每个阴性圆外区域立即取一个凸包。令$G_q$为严格位于$B(q,1000)$内的正32边形，用其各边依次分割当前位置多边形：每一步保存边外部分，再对边内部分继续分割，最后丢弃落在$G_q$内的部分。所得有限凸多边形并外包真实圆外区域，同时保留分开的候选位置及大部分孔洞。内接近似会保留圆弧附近的薄层，仍是保守处理，并非精确圆弧布尔运算。

光学失败同样加入原始历史，以20米排除圆进行上述分割。覆盖和行动使用19.8米半径留量；实际清除可能在19.8至20米之间成功，所以后续全局位置界采用$1800+20=1820$米，不能把全部实际成功位置都限制在1819.8米内。

\subsection{位置与发射方向的联合分区}
以有限并而非独立乘积表示混合源的相容状态：
\begin{equation}\label{eq:jointposterior}
\widehat H=\bigcup_{j=1}^{m}\left(\widehat P_j\times I_j\right),
\qquad I_j=[\alpha_j,\beta_j].
\end{equation}
每个方向区间具有自己对应的位置多边形，全向假设另作无方向分支。普通阴性$q$被分为两个可能原因：超出同一固定接收半径，或位于定向不可见侧。超距分支应用式\eqref{eq:exactomni}中的1000米圆外及全部阳性距离比较；不可见分支仅保留定向假设。两者取并，不能在尚未排除超距时强行使用方向阴性。

对于某分区，若$q$到全部位置均小于1000米，或存在旧阳性$p$使全部位置满足$\norm{x-q}<\norm{x-p}$，则$q$必在接收距离内。后一条件平方相减后为线性式，在多边形顶点检查即可。此时阴性只保留不可见分支，全向分支可被排除。

令区间中点方向为$n_0=u((\alpha+\beta)/2)$，并定义
\begin{equation}
\eta_I=2\sin((\beta-\alpha)/4),\qquad
M_P(q)=\max_{v\in V(\widehat P)}\norm{q-v}.
\end{equation}
由$\norm{n(\theta)-n_0}\le\eta_I$，得到保守线性收缩：
\begin{align}
p\text{阳性}&\ \Longrightarrow\ n_0^\mathsf T x\le n_0^\mathsf T p+\eta_I M_P(p),\\
q\text{距离内阴性}&\ \Longrightarrow\ n_0^\mathsf T x\ge n_0^\mathsf T q-\eta_I M_P(q).
\end{align}
区间越窄，方向不确定带来的放宽越小。分割方向区间后，对各子区间重新应用历史约束，因此不同方向对应的位置区域能够逐渐分离。

第七版把细化预算与即将执行的动作连接。每个真实观测版本最多使用8次角二分，并维持48分区上限；常规观测更新仅重用约束，不自动按最大角宽细分。对下一搜索点、后续目标及局部测点，优先选择可能同时产生可见和不可见预测的分区，或达到最坏两段移动距离支撑值的分区。每轮只试分优先级最高的4个分区，用子分区的最坏移动支撑值与可见性放宽量是否下降来选择一次实际细分；候选行动最多4点，无可确认下降即停止。

这里的评分用于分配计算量，不作为概率、测量收益或真实成本下降证据；决定是否丢弃位置的仍是历史观测与保守几何约束。若某次分裂超过分区配额，保留该父分区和全部原始观测，不丢弃可能分支凑数量；后续新阳性到来时继续重用旧阴性。该策略有明确的计算规模，但配额耗尽仍会留下部分保守区域，不能声称完整精确恢复联合相容集。

为了复用单次清除和锚点收紧，可以暂取位置并集的凸包$\widehat P$。这不会额外削弱“一个圆覆盖全部可能源”的判据：对任意固定圆心，凸包上的最大距离与原并集相同。保留非凸分区的主要价值在于多圆终局、观测条件分支及位置方向联合排除，而非声称凸包必然改变每一个单次鲁棒清除解。

当前锚点仍可按$L\leftarrow\min\{L,\max_{v\in V(\widehat P)}\norm{v-a}+10^{-7}\}$收紧。进入连续局部定位后不再插入共享测向；此前的共享阳性可建立一个距离上界不超过1000米的新锚点，并保留全部历史。它没有重置任何已经消耗的局部收缩轮数。

\subsection{从估计源中心改为选择清除动作}
对历史区域所有可能源都可清除的位置集合为
\begin{equation}\label{eq:K}
K(\widehat P)=\bigcap_{v\in V(\widehat P)}B(v,r_s).
\end{equation}
由于距离函数为凸函数，覆盖全部顶点等价于覆盖整个多边形。该条件对外包区域精确，对真实相容集合保守。若交集非空，求解
\begin{equation}\label{eq:projection}
z^*=\arg\min_{z\in K(\widehat P)}\frac12\norm{z-p}^2.
\end{equation}
这是闭凸集上的欧氏投影\cite{boyd}，目标严格凸，非空时最优点唯一。当前位置已可行则原地清除；否则最优点在一条有效圆弧内部或至少两条圆弧交点上。因此可以枚举当前点到单个圆盘的投影以及两圆交点，再检验所有顶点距离，取最近可行候选。

对于等半径圆心$v_i,v_j$，令$d=\norm{v_j-v_i}$。在$0<d\le2r_s$时交点为
\begin{equation}
z_\pm=\frac{v_i+v_j}{2}\pm
\sqrt{r_s^2-d^2/4}\frac{(v_j-v_i)_\perp}{d}.
\end{equation}
若某点只在一条有效边界上达到最优，局部最优条件要求其为当前位置对该圆盘的最近投影；多个不同有效圆同时成立则属于上述交点。此枚举覆盖最优解，而不仅是任意候选采样。$m$个顶点的候选数为$O(m^2)$，可行性检验每次$O(m)$，故总复杂度$O(m^3)$。

保留式\eqref{eq:terminal}给出的独立候选。第七版同时计算历史区域面向当前位置和后续目标的可行投影，与锚点候选比较两段路径；执行与成本上界中选中的同一个点。历史交集暂不可用时继续确定收缩；历史区域若为空则视为模型或数值异常，不丢弃历史后强行执行。



\subsection{连续清除位置与任务调度的接口}
已能确定清除的源以最近鲁棒可行点及面向后续目标的可行点作为有限服务模式；尚未完成定位的源保留有保证的第一测点与其连续完成成本。旧有固定邻点下的凸下降和批次路径工具作为历史基准保留，第七版默认执行下文的完整任务图。

统一图并不把源位置估计中心当作清除点。每种模式必须给出实际入口、条件成本上界和包含所有成功结束位置的出口区域；只有这些可执行服务才参与任务顺序比较。

\section{问题三：七点搜索与未知频道排除}
\subsection{更短的七点构造及条件最小性}
采用
\begin{equation}\label{eq:q3}
Q_3=\{(0,0)\}\cup\{1125u(k\pi/3):k=0,\ldots,5\}.
\end{equation}
原点覆盖$r\le1000$。其余源与最近外环点的夹角不超过$\pi/6$，距离平方不超过
\begin{equation}
f(r)=r^2+\rho^2-2r\rho\cos(\pi/6),\qquad \rho=1125.
\end{equation}
$f$在$r\in[1000,1800]$上的最大值位于端点，最大距离小于999.111米，故七点足以发现全部全向源。使圆域边界被六点覆盖的最小环半径为
\begin{equation}
\rho_* = 1800\cos(\pi/6)
-\sqrt{1000^2-1800^2\sin^2(\pi/6)}
\approx1122.956.
\end{equation}
取1125保留接收裕量。从原点沿环依次访问六点的开放路径长$6\rho=6750$米，较旧1200米环的同构路径短450米，对应90秒纯移动时间。

还可得到一个有范围限定的最小性结论：若静态检测方案包含原点，半径1000米的其他检测圆在半径1800米的边界圆上至多覆盖$2\arcsin(1000/1800)$弧度。五个这样的圆总覆盖角小于$2\pi$，原点检测又不覆盖边界，因此至少需要原点加六个非原点检测位置。式\eqref{eq:q3}达到该点数下界。此结论不是任意起点、任意自适应路径或总时间的最优性结论。

\subsection{未知频道的连续区域判据}
未知频道在全部已扫描位置均无信号。对全向源，只要这些位置的1000米接收圆覆盖$D$，就可排除该频道，无须将“发现了10个源”当作停止条件。所有仍未知频道有相同的全局扫描历史，故可共用同一次覆盖证明。

程序把包围$D$的正方形递归四分。完全在$D$外的方框直接丢弃；若某个检测点到方框四角均小于1000米，则范数凸性保证整个方框均被覆盖。不能证明的方框继续细分；达到节点数或深度上限时返回“未决”，继续固定搜索。程序证明的是连续方框覆盖，没有用有限网格点全为阴性替代整个圆域的结论。

\section{问题四：25点环绕覆盖与搜索点替换}
\subsection{任意定向下的局部凸包判据}
对于位置$x$，令$Q_x=\{q\in Q:\norm{q-x}<1000\}$。若$x\in\operatorname{conv}Q_x$，则任意闭半平面方向至少包含一个检测点；否则相应位移向量在该方向上的内积均为负，其凸组合不可能为零。严格内点条件还可保证存在非重合且内积为正的检测点。

在要求严格接收裕量的设计中，这给出明确的几何目标：既覆盖距离，又使邻近测点在方向上包围源。它比普通圆盘并覆盖强，不依赖事先估计每个源的辐射方向。

\subsection{两环三角网构造}
令
\begin{equation}\label{eq:q4}
O=(0,0),\quad A_i=990u(i\pi/4),\quad
B_j=1836u(j\pi/8),\quad Q_4=\{O,A_i,B_j\},
\end{equation}
其中$i=0,\ldots,7$、$j=0,\ldots,15$。按下标循环，对每个$i$组成四个三角形
\begin{equation}
OA_iA_{i+1},\quad A_iB_{2i}B_{2i+1},\quad
A_iB_{2i+1}A_{i+1},\quad A_{i+1}B_{2i+1}B_{2i+2}.
\end{equation}
32个三角形无重叠地铺满外环正十六边形。外多边形内切圆半径为
$1836\cos(\pi/16)=1800.721775>1800$，所以整个源域严格包含其中。
三角网各类边长分别为
\begin{equation}
\begin{gathered}
990,\quad 2(990)\sin(\pi/8),\quad 2(1836)\sin(\pi/16),\quad846,\\
\sqrt{990^2+1836^2-2(990)(1836)\cos(\pi/8)}.
\end{gathered}
\end{equation}
最大值为996.210426米，严格小于1000米。

\begin{theorem}[25点混合源覆盖]\label{thm:25}
对源圆域中任意位置及任意闭半平面辐射方向，式\eqref{eq:q4}存在有效检测点，因此不依赖源总数即可完成未知频道搜索。
\end{theorem}
\begin{proof}
源所在三角形的各顶点到源距离均不超过该三角形最长边。源为顶点凸组合，故至少一个顶点位于可见半平面。若源恰为内部网格顶点，邻接三角形构成完整邻域，其非重合邻点也在1000米内并环绕源。源域严格在外多边形内部，故无需处理外边界顶点的缺失邻域情形。
\end{proof}

\begin{figure}[htbp]\centering
\begin{tikzpicture}[x=.00165cm,y=.00165cm]
\coordinate (O) at (0,0);
\foreach \i in {0,...,7}{
\coordinate (A\i) at ({990*cos(45*\i)},{990*sin(45*\i)});}
\foreach \j in {0,...,15}{
\coordinate (B\j) at ({1836*cos(22.5*\j)},{1836*sin(22.5*\j)});}
\foreach \i in {0,...,7}{
\pgfmathtruncatemacro{\k}{mod(\i+1,8)}
\pgfmathtruncatemacro{\e}{2*\i}
\pgfmathtruncatemacro{\m}{2*\i+1}
\pgfmathtruncatemacro{\f}{mod(2*\i+2,16)}
\draw[gray!65] (O)--(A\i)--(A\k) (A\i)--(B\e)--(B\m)--(B\f)--(A\k)
 (A\i)--(B\m)--(A\k);
}
\draw[thick] (0,0) circle (1800);
\fill[blue] (O) circle (25);
\foreach \i in {0,...,7}{\fill[blue] (A\i) circle (25);}
\foreach \j in {0,...,15}{\fill[blue] (B\j) circle (25);}
\node[anchor=west] at (2100,1100) {内环：8点，半径990米};
\node[anchor=west] at (2100,650) {外环：16点，半径1836米};
\node[anchor=west] at (2100,200) {实线圆：源分布区域};
\node[anchor=west] at (2100,-250) {32个三角形覆盖全部圆域};
\node[anchor=west] at (2100,-700) {最长边：996.211米};
\end{tikzpicture}
\caption{25点构造：三角形短边保证距离，凸组合保证任意方向可见}
\end{figure}

一条明确可行的开放访问顺序为$O,A_0,\ldots,A_7,B_{14},B_{15},B_0,\ldots,B_{13}$，长度为
\begin{equation}\label{eq:searchlength}
\ell_{25}=990+7\cdot2(990)\sin(\pi/8)+846
+15\cdot2(1836)\sin(\pi/16)=17885.5673.
\end{equation}
旧31点三角晶格任意两点间距至少990米；从原点出发访问其余30点的任意路线长至少29700米。因此新显式路径相对该下界短11814.4327米，即39.78\%。这证明搜索点集和一条访问路径的改进，不证明25点数目最少，也不证明穿插定位清除后的整局路程改善39.78\%。

\subsection{剩余排除区域驱动的搜索计划改写}
对问题四的一个方框，只选择到四角均小于1000米的共同邻近测点。若它们的凸包严格包含四角，则包含整个方框；其中任意源均满足局部凸包判据。提前结束仍使用最多2048方框、深度13的完整区域证明。

收尾计划不再只在清除后尝试原地换点。将未知频道的共同阴性历史用于四分区域，最多处理512个节点、深度8；已证明排除的整框移除，所有未决叶框以及预算耗尽时未处理的栈框完整保留。所得并集$\widehat Z$外包尚可能存在源的区域。它可能包含已实际排除但尚未证明的空间，未决框中心不能被当作源见证，面积也不是漏源概率。

候选新点取当前位置、当前位置与被替换点的中点及最大未决框中心，并限制在全局位置界内。替换点在全部有限插入位置中取路线最短者。删除与替换两类分别保留最有希望且节约超过1米的方案，先给替换一个证明名额，再给删除一个名额，避免大量便宜但不可证的删除提案挤占替换机会。

单次证明最多4096方框、一次决策最多两次调用。设已访问$v$批、初始扫描配额为$Q$、已用证明次数为$n$，本轮额度为
\[
\min\{2,[\min(16,4(1+\lfloor4v/Q\rfloor))-n]_+\}.
\]
因而证明计算随搜索进度释放，前期未通过不会耗尽后期机会；整局仍至多16次调用。只有出现新的真实扫描或清除，才允许再次提案。

只有“已访问点与修改后的全部剩余点仍覆盖整个源域”被证明时才采用。证据不足就维持旧计划；删点减配额，换点一对一消耗旧配额，均不增加扫描动作。每次改写保持全域覆盖，因此即使所有提前停止都未决，最终计划仍足以排除未知频道。局部路线节约单列为计划量，实际尾部排除成本由末次成功清除后的真实虚拟时间差计算。

\section{完整算法、实现与理论性能}

\subsection{对整个相容区域保证收缩的最近测点}
经典收缩只利用锚点距离上界$L$。当多条历史观测已显著压缩位置范围时，继续返回旧锚点附近会浪费信息。令位置外包多边形为$P$，$T(L)=0.502L-2.4$。对$r\le T(L)$定义
\begin{equation}\label{eq:regionalprobe}
K_r(P)=\bigcap_{v\in{\rm vert}(P)}B(v,r),\qquad
q_r=\arg\min_{q\in K_r(P)}\norm{p-q}.
\end{equation}
范数凸性保证$q_r$到整个$P$的距离不超过$r$，并非只保证对有限源样本成立。将半径缩放到已有清除圆交投影问题，可以枚举有限圆弧边界候选求最近点。实际计算内缩半径并保留正余量。

程序试用$r\in\{44.5,89.9,179.9,T(L)\}\cap(5,T(L)]$。问题三因$r<1000$保证接收；问题四还要求每个联合分区中不可见半平面与位置区域的交为空，从而全部相容发射方向均可见。无法证明时保留原成对测点。

\begin{proposition}[区域测点不增加核心测量上界]
对满足上述距离与可见性条件的$q_r$，一次测量必为近距离或正常响应。近距离直接同点清除；正常响应以实际示向和$L'\le T(L)$建立新锚点，同时已消耗轮数加一。因此该替代至少达到原阳性收缩速度，不增加五次全向、十二次混合核心测量上界。
\end{proposition}
该方法与增加一次试探不同：它直接替换本来必须执行的一次核心测量。候选半径仍是有限选择，尚未求全部连续测点及未来观测树的全局最优解。

\subsection{全部未完成任务的有向服务图}
设$\mathcal J$包含每个已发现未清除源和每个仍须访问的搜索点。同一次比较中所有候选顺序均必须完成同一$\mathcal J$。第六版一个源与两批扫描组成的固定片段被撤去，改为每轮执行一个任务后重新建图。

任务$i$的一种模式$m$由入口$e_{im}$、从入口开始的服务时间上界$s_{im}$和结束位置外包$E_{im}$组成。扫描模式有$s_{im}=11|U|$、$E_{im}=\{e_{im}\}$，覆盖当前每个未知频道最多6秒测向及5秒同点近距离清除。正常新发现源的后续处理成为下一轮的新任务，其成本不冒充已纳入当前图。

已能保证清除的模式取$s_{im}=5$和单点出口。测向模式直接在第一个实际测点进入，$s_{im}$按近距离、阳性及定向阴性后的第二测点分支取最大值，然后接连续仿射完成界；无需先虚构一次返回发现锚点的移动。普通和区域测点模式均用$E_{im}=P_i\oplus B(0,20)$包含全部可能成功清除位置。

模式间有向边权和起始权为
\begin{align}
w_{im,jn}&=\max_{z\in E_{im}}\norm{z-e_{jn}}/5+s_{jn},\\
w_{0,im}&=\norm{p-e_{im}}/5+s_{im}.
\end{align}
若$E=P\oplus B(0,20)$，最大距离恰为$\max_{v\in{\rm vert}(P)}\norm{v-e}+20$。这避免用估计中心制造过于乐观的源间路程。固定任务顺序后，沿模式层作动态规划，精确求该顺序的最小上界，终点自由。

\paragraph{同一有限图内的可计算间隙。}
对任务间边取全部模式对的最小权，引入起点虚拟节点及零成本返回边，得到最小指派松弛。松弛允许不连通环，也允许同一任务进出边采用不一致模式，故其值$\underline V_{\rm graph}$不超过完整可行计划的最小上界。以交换后继边连接不同环，构成可行种子，再与完整旧计划、贪心顺序、先源后扫描和先扫描后源比较。最多三轮任务搬移，每轮精算至多16个有希望的候选，最终由完整模式动态规划决定接受。

旧计划先补齐新发现任务再比较；历史证据只缩小时保留其仍有效的旧模式，不重置已消耗轮数。于是得到
\[
\underline V_{\rm graph}\le V_{\rm graph}^*\le
\overline V_{\rm selected}\le\overline V_{\rm incumbent}.
\]
左右两端是\emph{当前有限保守服务图}的优化界，不是实际整局最优虚拟时间的上下界。模式离散、出口外包与未来尚未发现任务都会影响与真实最优路线的差距。

\begin{figure}[htbp]\centering
\begin{tikzpicture}[>=Stealth, every node/.style={font=\small},
 box/.style={draw,rounded corners,align=center,text width=3.55cm,minimum height=1.05cm}]
\node[box] (a) at (0,0) {全部待处理源\\与剩余扫描};
\node[box] (b) at (4.5,0) {入口、服务上界\\与整个出口区域};
\node[box] (c) at (9,0) {指派下界与模式规划\\选择完整任务顺序};
\node[box] (d) at (4.5,-2) {执行首个任务\\记录成本并更新证据};
\draw[->] (a)--(b);\draw[->] (b)--(c);
\draw[->] (c.south)|-(d.east);\draw[->] (d.west)-|(a.south);
\end{tikzpicture}
\caption{第七版按单个完整任务滚动执行；未发现源在真实出现后入图。}
\end{figure}

\paragraph{执行与预测的双重约束。}
选中首模式后分别记录独立清除上界$U_0$和清除后到下一计划入口$t$的上界$U_t$。实际执行只完成首任务，保存其真实耗时$C$与结束位置$z$，检查
\[
C\le U_0,\qquad C+\norm{z-t}/5\le U_t.
\]
第二式中的最后一段是从真实结束位置出发的\emph{可行后续路程}，并未声称已经到达$t$。下一轮可重排全部任务，从而允许多个源连续清除，也允许在整图成本更低时延后某源。日志保留发现至开始处理的等待时间，便于识别是否仍有长期积压。

\subsection{在完整响应范围内评价额外信息测点}
额外观测只发生在连续收缩开始前，之后立即连续完成该源，不中途重启轮数。除当前位置外，候选包括最近保证接收区域投影、窄侧移点及一个较宽交会点，均限制在全局位置界内。已在同一点取得固定示向的观测不重复。

接收保证既可由所有候选位置均在1000米内给出，也可利用同一源固定半径：若历史阳性点为$a$，且对整个候选区域均有$\norm{x-q}\le\norm{x-a}$，则在$q$也不因距离失去接收。平方距离相减为仿射函数，整个凸分区可由顶点判定。定向不可见仍保留为阴性分支。

正常输出划为覆盖完整允许示向范围的连续区间，初始至多4段，优先细分当前最坏正常分支至最多24段。对每段$[\ell,h]$用扩大楔形$[\ell-\varepsilon,h+\varepsilon]$形成外包后验，同时保留可能的近距离与定向阴性。区间中点不能当作实际示向证明清除圆。

正常分支仅以新点和不超过1500米的距离上界评价角度无关的仿射完成策略；阴性沿用旧锚点，近距离直接清除。每个分支的独立成本和带目标成本来自\emph{同一条可执行延续}，不能分别取两个不同方案的最小值。对两类成本分别取全部分支最大值并加入移动、测量和切换成本；只有两者均较当前所选服务模式改进超过1秒才增加观测。执行中直接清除、区域测点和光学替代也必须同时满足这两个预算。

该严格门槛可能不触发。前瞻未触发应解释为当前外包界不足以确认收益，而非无须信息或策略已达最优。每源两次额外观测配额为保守上界；当前默认一次额外观测后立即完成该源。

\subsection{非凸区域上的有限光学覆盖搜索}
第七版保留狭长矩形的解析圆链：短轴半宽$w<r_s$，令$h=\sqrt{r_s^2-w^2}$，长度$d$的矩形可由$\lceil d/(2h)\rceil$个适当排布的等半径圆覆盖，最多允许三圆。除此之外，直接使用联合后验中的分开位置区域，避免为了覆盖两个相隔较远的小区域而把中间空白也当作必须覆盖的目标。

程序沿位置区域的主轴把凸分区再分成有限个凸单元，从各单元的鲁棒清除位置及矩形圆链生成候选，保留至多14个候选清除点。在最多三步的候选顺序中作有限搜索。只有当每个连续凸单元都被某个候选圆完整覆盖时才接受；范数凸性使“完整覆盖一个凸单元”可由其全部顶点验证。这里顶点用于证明单元覆盖，不是将若干源位置样本全部覆盖冒充连续区域覆盖。

第七版补充主轴及锚点方向上的二、三条连续切片覆盖；旧有14点候选和矩形方案始终保留。对成本较小的至多4个覆盖方案，进一步优化连续圆心。先按原圆心的Voronoi半平面把每个位置多边形完整分配给最近圆心，再将分配部分的全部顶点约束为该圆心必须覆盖的集合。每个圆心的可行域是有限圆盘的凸交，因此可在最近可行点、相邻两段路径下降点及它们与旧点的凸组合之间进行至多3轮坐标改进。

每步比较完整的最坏首次成功成本，而不只比较相邻两条边；只有该目标下降且分配区域仍全部覆盖才移动圆心。分区未通过覆盖余量检查或改进无效时保留原方案。连续改进扩大了可选位置，但有限切片、顺序和坐标迭代仍不等同于连续全局最优。

固定圆心序列$q_1,\ldots,q_k$，首次成功在第$j$步时的成本为
\begin{equation}\label{eq:opticalcost}
J_j=\frac15\sum_{i=1}^{j}\norm{q_i-q_{i-1}}+3(j-1)+5
+\frac15\norm{t-q_j},\qquad q_0=p.
\end{equation}
无后续目标时删去末项。第七版引入此前失败条件：令$P_0=P$，$P_j=P_{j-1}\setminus B(q_j,20)$。只有$P_{j-1}$仍可能在$q_j$成功时，$J_j$才是可能的首次成功分支；若整残余区域距$q_j$大于20米，该次尝试不可能成功，直接省去其动作与绕路。残余区域以保守多边形并表示，容量达到上限时保留父区域。

对原有连续覆盖方案的至多六种圆心排列逐一计算条件首次成功上界，并选择满足门槛的最佳顺序。省略确定无用的尝试不会破坏整个区域的原覆盖。有限候选与圆心改进仍不等同于所有连续圆心和全部决策树的最优方案。

光学失败后，20米排除圆立即进入联合后验。在剩余尝试预算内重新寻找覆盖；只有带后续目标和不带目标的两项新剩余上界均不超过原承诺上界减去实际已付前缀时才替换，否则执行原序列。每源总光学尝试数始终不超过三次。

\subsection{用测向策略下界控制光学切换}
仅有$U_{\rm opt}<U_{\rm RF}$不足以证明光学更好，因为测向上界可能很松。第七版改用下界门槛。先从相容区域中选取少量候选位置，对每个候选逐条检查源域、正常响应的5米下界、全部示向约束、光学失败距离、同点固定输出及同一个固定接收半径。令
$R_x=\max\{1000,\max_{p_i\in P^+}\norm{x-p_i}\}$；若某阴性不能由超出$R_x$解释，还须找到一个同时满足所有阳性可见和这些阴性不可见的方向。只有通过完整物理条件的候选才作为可行世界见证。外包顶点未经这些条件确认不能用来抬高下界。

对任意见证位置$x$，任何策略成功清除时都必须到达$B(x,20)$，因此其移动路线长度至少为
\begin{equation}
b_x(p,t)=\max\{\norm{p-t},\,
[\norm{p-x}-20]_+ +[\norm{t-x}-20]_+\}.
\end{equation}
没有后续目标时取$b_x(p)=[\norm{p-x}-20]_+$。设当前射频频道与目标频道不同的指示量为$s$，则对所有“先对本源测向，再完成清除”的单源连续策略，有最坏成本下界
\begin{equation}\label{eq:rf_lower}
L_{\rm RF}(H,p,t)=10+s+\frac15\max_{x\in W}b_x(p,t).
\end{equation}
其中10秒为至少一次测向及一次成功清除。未找到见证时只保留直接路径长度产生的通用下界，不用未确认的候选位置补足。有限见证只强化下界，不用于证明覆盖。

\begin{proposition}[条件最坏成本的光学优势门槛]
若一个连续区域光学覆盖满足$U_{\rm opt}(p,t)<L_{\rm RF}(H,p,t)$，则其最坏成本严格低于所有同源测向优先、且清除后到达同一目标$t$的策略的最优最坏成本。
\end{proposition}
\begin{proof}
见证世界均符合历史条件。任意测向优先策略在每个见证世界中均需支付上述动作成本，并经过源的20米邻域，故该策略的最坏成本不小于式\eqref{eq:rf_lower}。光学计划覆盖整个相容区域，其任意响应序列成本不超过$U_{\rm opt}$。上下界严格分离即得结论。
\end{proof}

\paragraph{圆邻域访问的对偶下界。}
原距离三角界会分别减去两次20米，忽略同一个清除位置必须同时服务进出两段路径。第七版对每个有效位置见证$x$，改用
\begin{equation}
\ell_x(p,t)=\min_{\norm{z-x}\le20}
\{\norm{p-z}+\norm{z-t}\}.
\end{equation}
不把有限数值优化得到的可行值误作下界。令$f(z)=\norm{p-z}+\norm{z-t}$，在任意试点取次梯度$g$，凸性给出
\begin{equation}
\ell_x(p,t)\ge f(z)+g^\mathsf T(x-z)-20\norm{g}.
\end{equation}
至多24次可行迭代中保留最大的上述对偶值，并与原距离下界取大。任意时刻停止均保留下界方向，浮点计算减去工程余量。新的测向策略下界为$10+c_s+\max_{x\in W}\ell_x/5$。例如$p=(-100,0),t=(100,0),x=(0,100)$时，原界为242.843米，圆邻域最短访问距离为$2\sqrt{100^2+80^2}=256.125$米，新对偶界能够收紧到该值附近。

见证生成也按路径成本选择空间极值并作有限内插扩展，每个候选仍逐条满足全部真实观测后才纳入；保留至多12个成本相关见证。多取一个有效见证只会收紧最坏成本下界，无法找到见证则保留通用下界，决不以不相容位置抬高门槛。

实际还同时要求去掉后续目标项时也满足该不等式，避免仅因后续路线变化就失去单源清除成本的比较依据。若上下界重叠，程序保留测向，不再因光学成本小于一个松的测向上界而切换。该门槛比较的是最坏成本，既不宣称每个真实位置上都更快，也不证明多源整局最优。

\subsection{有限完成与动作上界}
\begin{theorem}[第七版的有限完成条件]
在题设观测模型、精确几何判断以及足够执行预算下，第七版可以完成全部源的清除。每次连续定位仍受定理\ref{thm:finite}约束，且最多进行三次光学清除；全局搜索批数不超过7或25，服务图每轮完成一个已知源或一批扫描，所有新源在发现后入图。
\end{theorem}
\begin{proof}
初始搜索计划已被覆盖定理保证；删减或替换先证明新计划覆盖，再减少或消耗一个旧配额。提前排除仅在连续区域证书成功时生效。额外信息阳性只在连续定位开始前建立距离界不超过1500米的锚点，尚未消耗局部轮数；进入连续定位后不重启轮数，区域测点达到至少同样的阳性收缩量，故五次全向或十二次混合核心测量上界保留；任何光学替代均有有限圆覆盖保证。

主循环每轮至少访问一个剩余搜索点、经覆盖证明删除一个剩余搜索点，或成功清除一个已发现源。量
$|\text{剩余搜索点}|+16-|C|$为非负整数并严格减小，故主循环最多23或41轮。搜索已完成、16源均已发现或所有未知频道已由证书排除后，仍将剩余已发现源全部清除。结束时还要求发现数在10至16之间，否则报告不相容。
\end{proof}

最保守地将每批扫描均按20个频道计、每个源均按两次共享测向及三次清除计，得到
\begin{align}
A_3&\le7(20)+16(2)+16(5)+16(3)=300,\\
A_4&\le25(20)+16(2)+16(12)+16(3)=772.
\end{align}
以上动作数不含进入、退出接口，不把同一请求的幂等重试当作新的物理动作。更紧的实际计数会扣去已发现频道的后续扫描，本文保留简单统一上界。允许每源最多两次光学失败，故$C_f\le32$。

\subsection{虚拟时间的保守上界}
仅执行连续仿射策略，从当前位置到锚点距离为$d_0$时，对递推距离界逐轮求和可得
\[
F_3\le d_0/5+337,\qquad F_4\le d_0/5+704.
\]
第七版改用适合实际第一测点入口的统一服务界。全局任务之间的位置为搜索点或成功清除点，半径不超过1820米或1836米；待处理源的发现锚点也在该范围内。额外信息点受到相同限制，且紧接着连续完成该源。

当$L\le1500$时，原成对测点距锚点小于754米，二点间距不超过60米。故从全局当前位置直接进入第一测点的时间至多
$(2\cdot1836+754)/5=885.2$秒。给第一测量6秒、第二点移动12秒、第二测量5秒，以及双阴性后最保守的返回锚点和连续完成界，得局部服务至多
\[
6+12+5+754/5+704=877.8\ {\rm s}.
\]
合计小于1800秒。区域测点到每个相容源不超过750.6米，入口半径不超过2550.6米，其一次测量加连续完成也受上述宽松界控制；确定清除模式更短。保留的历史模式仍满足相同几何界。

信息替代同时降低所选模式的独立完成上界；光学上界低于测向策略下界，因而低于所选测向模式上界。连续处理中的替换按“已付前缀加剩余界”核算，故每个单源服务仍不超过1800秒。搜索批扫描按每频道11秒放宽，包含可能的同点近距离清除。对全部16源再计服务会重复计算部分近距离源，但只会使界更保守。因此
\begin{align}
T_3&\le 7(3640/5+20\cdot11)+16\cdot1800=35436,\\
T_4&\le 25(3672/5+20\cdot11)+16\cdot1800=52660.
\end{align}
两界均低于360000秒虚拟预算。它们依赖题设、几何条件和正确执行成本模型，不能预测真实成绩，也不保证任意网络条件下满足真实运行时限。

\subsection{数值、通信与结果记录}
几何证明使用精确算术；代码采用方向半宽$1.005^\circ$、外扩半平面、内缩清除半径及严格正余量。连续证书的浮点判断属于工程余量，尚非区间算术认证。相容区域为空、保证接收的全向测点无信号、保证清除失败等均报告模型不相容。

通信保持串行。响应不确定时，只允许相同请求标识与内容重试；未解决的请求禁止后续新动作。每次请求与响应、元数据及最终汇总分别存入动作日志、元数据文件和汇总文件。程序在真实时限结束前预留25秒，并逐动作检查预算；任意网络延迟下不可能由虚拟时间界推出真实20分钟内必定完成。

第七版记录完整已知任务的计划界、指派松弛下界、选中首模式、每次服务的实际成本与两项余量，并记录源等待时间、区域测点使用次数、额外信息位置和覆盖证明释放进度；联合分区、光学门槛及失败重规划记录继续保留。这些记录用于解释后续实测的收益来源；不将某次局部计划路程下降直接登记为整局实现收益。

\section{解析结果、必要检查与真实演练}
\subsection{解析比较与必要离线检查}
问题三七点的数目结论限于包含原点检测的静态全向覆盖；问题四25点只证明可行。问题二的886.3434米结果是解析计算；任何解析界下降均不能替代真实运行成绩。

旧版检查保留已有记录，第七版新增12项必要固定输入检查通过：小规模完整枚举核对指派下界及模式动态规划、多源连续处理、全区域测点与定向可见性、固定半径耦合、区域正常响应、定向阴性后近距离、完整示向区间及额外观测配额、光学条件分支与门槛、覆盖名额与阶段预算、问题二解析间隙、正式表版本隔离及预置响应下的完整服务记录。检查不生成隐藏物理环境，不连接模拟器。

\subsection{第三、六版真实演练与性能瓶颈}
用户提供的第三版及第六版运行日志分别独立保留。成功清除数来自明确成功响应，虚拟时间取最终值，$\overline T=T_v/N_c$；本地进程耗时不能替代官方程序运行时间。
\begin{table}[htbp]\centering\small
\caption{真实演练数据：不同案例之间不构成版本提速对照}\label{tab:practice}
\begin{tabular}{llrrrr}\toprule
版本&题号&清除数&虚拟时间/秒&移动距离/米&平均时间/秒\\\midrule
第三版&三&12&4046.411&16522.057&337.201\\
第三版&四&10&10933.864&42239.320&1093.386\\
第六版&三&13&4889.664&20798.319&376.128\\
第六版&四&14&8753.575&36667.873&625.255\\\bottomrule
\end{tabular}
\end{table}

第三版问题四有操作者截图确认总数10，因此支持该例全部清除。其余三例缺独立官方总源数，清除数与程序完成证书不能被当作已独立确认的100\%清除率。

第六版案例为问题三NF56-ST4Z-8BKW-B84E、问题四2N7V-DK3C-82EN-XS8C。两次均正常退出并给出完成证书，分别完成7、25批搜索；测量114、233次，清除尝试14、15次，各有一次失败后成功的光学尝试。移动时间占85.07\%、83.78\%，主要成本集中在路线与源处理移动。

问题三末批扫描时仍有10个已知源未完成，随后支出2546.523秒；问题四相应为2源和793.683秒。两例末次成功清除都发生在结束时，因此“末次清除后排除时间”为零，不能据此推断已消除无效搜索成本。第六版额外共享观测与覆盖改写均未采用，问题四16次覆盖改写证明配额已经耗尽。这些事实促使本轮改写调度、核心测点及证明预算。

\subsection{真实观测状态上的计划界分析}
只读取两份真实日志的第一批搜索结束状态，截断虚拟时间均为119秒；问题三已知7源、剩余6批搜索，问题四已知6源、剩余24批。把这些观测输入新服务图，比较相同已知任务集合的不同顺序。运输对象明确禁止执行任何动作，原响应不移植到新坐标。
\begin{table}[htbp]\centering\small
\caption{真实首批证据上的有限服务图成本界，非新版运行成绩}
\begin{tabular}{rrrrr}\toprule
题号&基准上界/秒&选中上界/秒&指派下界/秒&相对下界间隙\\\midrule
3&5882.595&4357.327&4306.499&1.180\%\\
4&12691.135&11437.960&11253.759&1.637\%\\
\bottomrule\end{tabular}\end{table}

表中基准为同一新服务图内“先扫描、再处理已知源”的完整顺序，不是第六版实测剩余耗时；选中界覆盖全部当前已知任务。指派间隙仅针对有限保守图的最优上界。问题三选中一个区域测点模式，两个首批状态均未接受额外信息观测；这种局部触发记录不代表整局使用率。

新发现源的服务尚未包括在首批图内，执行后重新发现和重规划的实际路径也未发生。因此不据此生成第七版整局时间、清除率或真实提速百分比。分析展示的是新规则能够比较并保留较低的同任务计划界。

\subsection{第七版实测接入}
人工双击入口默认策略标识为\texttt{obligation-service-interception-v7}。操作者选择模拟器模块、输入案例编码并确认连接，本次开发没有启动这些入口。新版汇总增加单源等待、区域测点、信息观测、任务图间隙以及每次服务的两项成本余量。

结果整理按策略版本与演练/正式类型隔离。官方总源数、官方程序运行时间和日志文件名从操作者独立记录取得，缺失保留待补。第六版演练已经接入，下一次需由用户完成第七版问题三、四各一次演练，再按题面要求提供正式结果；不同案例不可直接计算版本提升比例。

\subsection{六次正式测试表：等待实际数据}
按题面要求分别报告三次正式测试。程序运行时间必须取官方记录，与式\eqref{eq:cost}的虚拟时间及本地进程耗时区分。正式测试不公开真实总源数，故不额外填报未知分母的清除率。以下表格尚不能用于最终提交。
% BEGIN_V7_FORMAL_RESULTS
\begin{table}[htbp]\centering\caption{问题三正式测试结果，尚未取得}
\begin{tabular}{lcccc}\toprule
序号&测试案例编码&清除干扰源个数&平均定位清除时间/秒&程序运行时间/秒\\\midrule
测试1&\pending&\pending&\pending&\pending\\
测试2&\pending&\pending&\pending&\pending\\
测试3&\pending&\pending&\pending&\pending\\\bottomrule
\end{tabular}
\end{table}
\begin{table}[htbp]\centering\caption{问题四正式测试结果，尚未取得}
\begin{tabular}{lcccc}\toprule
序号&测试案例编码&清除干扰源个数&平均定位清除时间/秒&程序运行时间/秒\\\midrule
测试1&\pending&\pending&\pending&\pending\\
测试2&\pending&\pending&\pending&\pending\\
测试3&\pending&\pending&\pending&\pending\\\bottomrule
\end{tabular}
\end{table}
% END_V7_FORMAL_RESULTS
六个官方日志须由操作者导出并保留原文件名，与案例编码一一对应。参考仓库记载的其他运行结果不属于本文算法，不能代填本表。


\section{模型讨论与适用边界}

\subsection{本轮贡献及其边界}
本轮把真实积压与移动成本对应到具体机制：所有已知任务共同规划；以真实首测点入图消除虚构返回；用整个相容区域选择最近保证收缩测点；完整响应范围内评价主动信息；按失败条件简化光学序列；为后期覆盖修改保留计算机会。问题二则把所需精度直接写成约束，给出明确的移动距离优化间隙。

凸投影、指派、动态规划、分支定界和路径局部改进均为已有方法。本文针对题设提出服务入口与不确定出口、轮数保持的区域探测及两项执行预算之间的具体连接，未完成全面文献排查，不能宣称这些组合已经得到学术首创确认。

\subsection{仍保留的重要不足}
第一，有限服务图最优不等于物理整局最优。模式离散、出口区域取最大距离、当前未知源未入图，会造成保守性和短视；指派下界也未处理模式连贯与子环限制。首批图可以仍先完成较长搜索段，撤销固定配额并不意味着必然消除实际积压。

第二，区域测点采用四个半径候选。定向可见性证明可能太强，未证明时退回成对探测；完整信息前瞻也可能因上下界较松而拒绝实际上有益的观测。当前并未求解任意观测树或同时服务多个频道的连续信息路线。

第三，48个联合分区、角区间放宽、内接排除圆及有限见证会影响效率。达到配额时保留父区域防止漏掉真实世界，却保留部分实际上已排除的空间；浮点工程余量仍非严格区间算术。

第四，光学只搜索至多三次覆盖方案，失败条件剪枝与连续圆心改进不等于全局最优决策树；25点最小性、整局近似比及精确物理直径极小极大尚未证明。最重要的经验缺口是尚无第七版真实演练，不能凭机制复杂程度推断平均成绩。

\subsection{适用条件与实测缺口}
闭半平面辐射、静止源、固定频道、接收半径下界和有界误差是当前保证的基础。一般径向区间$[r_0,L]$可用$A=(L+r_0)/2$、$B=\beta L$，但必须重新满足
\begin{equation}
\max_{r\in\{r_0,L\}}\{(r-A)^2+B^2+2Ar(1-\cos\varepsilon)
+2Br\sin\varepsilon\}<R_{\min}^2,\quad B>A\tan\varepsilon.
\end{equation}
窄波束、障碍物或时变源不能直接使用本文次数定理。浮点实现尚缺严格区间认证，通信异常也可能导致保守停止。

已有四次真实演练分别属于第三、六版，尚缺第七版演练及题面规定的六次正式结果。论文因而能够陈述实现和条件保证，暂不能量化新版的实际优势。下一项最有价值的实测是由用户运行第七版并提供原始记录，先判断总时间、移动时间及尾部排除成本是否按预期下降，再决定是否增加必要的对照。


\section{结论}
本文形成第七版搜索与清除算法：以连续覆盖证书保证搜索，以完整已知任务服务图协调顺序，以全区域最近保证收缩测点利用累积历史，以完整响应区间和双成本约束控制额外观测，并用条件首次成功分支改进光学终局。有限服务图的可计算间隙与每次任务的真实成本分别记录，避免把计划下降当作已经实现的收益。

问题二在163米解析精度目标下给出886.3434米设计及小于0.0820米的计算间隙。问题三、四保留300、772次动作上界，并得到35436、52660秒的保守虚拟时间界。12项新版必要检查通过；第三、六版真实演练已接入，但第七版实际优势仍须由操作者运行确认。本文没有证明整局全局最优，也没有把尚未进行的正式测试写成已取得成绩。

\begingroup\small
\begin{thebibliography}{9}\setlength{\itemsep}{0pt}
\bibitem{problem} 全国大学生数学建模竞赛组委会. 2026年高教社杯全国大学生数学建模竞赛B题：无线电干扰源的快速自动定位与清除[Z]. 2026. 用户提供题面PDF.
\bibitem{manual} 2026年全国大学生数学建模竞赛B题附件1：模拟器使用说明[Z]. 2026. 用户提供附件.
\bibitem{protocol} 2026年全国大学生数学建模竞赛B题附件2：HTTP+JSON通信协议[Z]. 2026. 用户提供附件.
\bibitem{boyd} Boyd S, Vandenberghe L. Convex Optimization[M]. Cambridge University Press, 2004. \url{https://web.stanford.edu/~boyd/cvxbook/}.
\bibitem{repo} Aze-Sh. MCM-：数学建模参考仓库[CP/OL]. 用户提供的本地下载版本，访问日期2026-09-11. \url{https://github.com/Aze-Sh/MCM-}. 参考其历史外包区域思路，本文策略与代码另行实现.
\bibitem{cetsp} Gutow G, Choset H. Efficient Second-Order Cone Programming for the Close Enough Traveling Salesman Problem[CP/OL]. CMU作者提供论文，访问日期2026-09-11. \href{https://biorobotics.ri.cmu.edu/papers/paperUploads/Efficient_SOCP_for_the_CETSP_ICRA_2025_v5_Geordan.pdf}{CMU作者提供的论文全文}.
\end{thebibliography}\endgroup

\clearpage
\appendix

\section{代码与结果字段对应}
\begin{longtable}{p{.28\textwidth}p{.63\textwidth}}\toprule
文件&功能\\\midrule
\texttt{geometry.py}&楔形交分类、凸包、直径与直径圆判断\\
\texttt{solve\_geometry.py}&问题一、二离线几何入口\\
\texttt{strategy.py}&保留基线与第三版搜索和仿射收缩原语\\
\texttt{coverage.py}&第七版7/25点设计、三角网、连续区域证书\\
\texttt{evidence.py}&圆域外包、阴性圆凸包、定向半圆证据\\
\texttt{posterior.py}&历史区域裁剪、距离收紧、最近鲁棒清除\\
\texttt{planning.py}&开放路线、源插入、光学覆盖和剩余成本界\\
\texttt{routing.py}&固定邻点下的有界凸下降\\
\texttt{cooperative.py}&保留第四版及共用覆盖、联合清除、完成流程\\
\texttt{joint\_belief.py}&非凸位置方向分区、历史重用及可行世界见证\\
\texttt{decision.py}&整类响应前瞻、有限光学覆盖与测向下界\\
\texttt{adaptive.py}&保留第五版执行流程\\
\texttt{adaptive\_v6.py}&保留第六版同终点片段\\
\texttt{tail\_search.py}&剩余区域、证明后删减与替换\\
\texttt{refinement.py}&行动相关几何细化预算\\
\texttt{decision\_v6.py}&圆邻域对偶下界与连续圆心改进\\
\texttt{q2\_design.py}&固定移动距离的角区间分支定界\\
\texttt{q2\_budget.py}&精度约束下移动距离分支定界\\
\texttt{service\_graph.py}&完整服务图、指派松弛与模式动态规划\\
\texttt{probe\_frontier.py}&全区域最近收缩测点及定向可见性\\
\texttt{interception.py}&主动信息候选及完整响应成本界\\
\texttt{optical\_v7.py}&光学失败条件及双预算重规划\\
\texttt{coverage\_v7.py}&覆盖提案类别名额与分阶段证明预算\\
\texttt{adaptive\_v7.py}&单任务滚动执行、服务成本与等待记录\\
\texttt{observed\_\allowbreak{}state\_\allowbreak{}analysis.py}&真实观测状态的只读计划界分析\\
\texttt{collect\_results.py}&只读真实日志及正式结果成表\\
\texttt{solver.py}&第三版执行流程及共用状态、动作记录\\
\texttt{protocol.py}&串行接口、相同请求重试及未决请求隔离\\
\texttt{run\_robot.py}&人工启动入口、策略选择、结果输出\\
\texttt{test\_offline.py}&保留的14项必要检查\\
\texttt{test\_v4.py}&保留第四版8项固定输入检查\\
\texttt{test\_v5.py}&保留第五版12项固定输入检查\\\bottomrule
\end{longtable}


默认策略为\texttt{obligation-service-interception-v7}。Python参数\texttt{--v6}保留第六版，\texttt{--v5}、\texttt{--v4}、\texttt{--v3}及\texttt{--baseline}分别运行更早版本，选项互斥；PowerShell对应参数以单横线及大写V表示。第六版源码和论文冻结在\texttt{versions/v6}，更早备份保留。

默认启动器与整理后的源码在同一完整项目文件夹，运行数据写入其“运行结果”，不写回其他旧工程。第七版12项检查、问题二解析区间和真实首批观测分析在“测试记录”下保留。汇总中的\texttt{service\_receipts}记录真实服务和可行后续段，\texttt{source\_waits}记录已发现源的等待；旧片段字段仅为历史兼容。

本地\texttt{local\_elapsed\_s}不替代官方程序运行时间。官方日志由操作者导出并保持原名；跨进程断点续跑尚未实现。

\section{AI辅助与参考材料使用说明}
本文建模、代码和文字整理使用了ChatGPT/Codex辅助。题面、两份附件及用户提供的参考仓库作为材料输入；仓库启发了历史外包区域的使用，其运行成绩未作为本文成绩。已有凸几何和路径优化工具保留来源，针对本题的设计与组合不等同于经全面检索确认的学术首创。

第三、六版共四次演练均由用户操作模拟器完成，本次第七版开发未启动或连接模拟器。新增检查为固定输入和预置接口响应，不生成虚构正式日志。正式提交前应按实际协作记录完成独立AI工具使用详情材料，并补入真实测试数据。本段不替代该独立材料。
\end{document}
